\section{规格编写工作量结果}
\label{sec:resultados-esfuerzo}

十个官方配置合计包含 4074 行非空且不含注释的代码（LOC），而
\texttt{.m2b} 模型为 1121 LOC。应用公式~\ref{eq:reduccion-loc} 后，
加权 LOC 缩减率为 72,48\,\%。各案例缩减率的平均值为 73,92\,\%，中位数为
75,40\,\%。按 UTF-8 字节数计算，缩减幅度相近：从 137358 字节降至
40563 字节，即缩减 70,47\,\%。

上述总数将每个编辑器视为可独立维护的产物。为避免 Pond Tutor 与 Pond
之间官方实现中的代码复用使 M2B 获得人为优势，研究剔除共享的官方代码行后重新进行了
计算。此时，保守计算的官方代码总量变为 3737 LOC，缩减率仍达到
70,00\,\%（按字节数计算为 67,67\,\%）。因此，这一趋势并不依赖于对
337 行共享代码进行两次计数。

这一结果需要结合功能一致率来理解。Model2Blockly 大幅缩减了需维护的源代码
规模，但这些模型并不能精确替代所有官方配置。该语料集同时表明，描述工作量
有所减少，并揭示了一组仍然缺失的具体能力；但它并不能证明，只需减少
72,48\,\% 的代码便可实现每个编辑器 100\,\% 的功能。

\section{语料集中观察到的局限}
\label{sec:limites-corpus}

表~\ref{tab:limites-editor} 汇总了编辑器中最常见的差异。这里的计数表示受影响的
原子属性数量，而不是相互独立的特性数量：同一项能力缺失可能会影响一个案例中的
多个块和字段。

\begin{table}[htbp]
  \centering
  \caption{造成主要差异的编辑器能力}
  \label{tab:limites-editor}
  \resizebox{\textwidth}{!}{%
  \begin{tabular}{llrr}
    \toprule
    能力 & 主要状态 & 属性 & 案例 \\
    \midrule
    特定连接策略 & \texttt{unsupported} & 101 & 9 \\
    类别高级组合 & \texttt{unsupported} & 66 & 5 \\
    显式虚拟输入 & \texttt{unsupported} & 52 & 4 \\
    只读调用标签 & \texttt{partial} & 30 & 2 \\
    有意省略类别颜色 & \texttt{unsupported} & 30 & 5 \\
    工具箱整体组合 & \texttt{unsupported} & 19 & 2 \\
    工具箱中的预配置块 & \texttt{unsupported} & 19 & 7 \\
    输入对齐 & \texttt{unsupported} & 14 & 3 \\
    阴影块配置 & \texttt{partial} & 11 & 2 \\
    命名虚拟输入 & \texttt{unsupported} & 11 & 3 \\
    保留块的重新定义 & \texttt{mismatch} & 8 & 3 \\
    \bottomrule
  \end{tabular}%
  }
\end{table}

影响范围最广的局限是连接策略。通用元模型能够表示一个块是否具有输出连接、
前序连接和后继连接，但无法表示官方编辑器通过 JavaScript 添加的所有动态检查与
例外规则。以编程方式构建的工具箱、带初始状态的阴影块、输入的视觉对齐方式，
以及仅用于控制布局的虚拟输入，也都难以处理。这些特性不会阻止编辑器加载，
但会妨碍完整的结构复现。

在生成器方面，最常出现的差异包括：空输入的默认值（六个案例中的 27 个属性）、
枚举映射（三个案例中的 10 个属性）、模型自动序列化（Puzzle 中的六个属性）、
受变形器状态影响的生成（Pond Tutor 和 Pond 中的六个属性），以及过程和变量数据库。
这些观察结果指出了可落实的改进方向：扩展输入与工具箱的元模型，引入声明式连接
策略，并提供有状态的生成服务。

\section{输入路线的等价性}
\label{sec:resultados-rutas-entrada}

RQ4（输入路线）比较的不是 Ecore 与 DSL 的文本规模，而是两个适配器的结果。
研究选择 Graph、Maze 和 Turtle，以逐步覆盖简单字段、连接、下拉菜单、带类别的
工具箱、初始工作区和生成器。在三个案例中，均严格比较了
表~\ref{tab:equivalencia-rutas} 所列规范描述符的七个部分。

\begin{table}[htbp]
  \centering
  \caption{Ecore 与 \texttt{.m2b} 输入路线之间的严格等价性}
  \label{tab:equivalencia-rutas}
  \begin{tabular}{llll}
    \toprule
    案例 & 复杂度 & 已比较部分 & 结果 \\
    \midrule
    E01 Graph  & 简单 & 7/7 & 等价 \\
    E03 Maze   & 中等 & 7/7 & 等价 \\
    E07 Turtle & 复杂 & 7/7 & 等价 \\
    \bottomrule
  \end{tabular}
\end{table}

结果表明，在 3/3 个案例中，两条输入路线均等价：运行控制参数、块、工具箱、初始工作区、
生成器、错误和可加载性
完全一致。该测试也有助于发现实现缺陷。在 Graph 中，研究发现 Ecore 模型缺少
一个 \texttt{output} 注解；在 Turtle 中，则修正了适配器对下拉菜单标签和值的
解释。修正这些问题后，研究重新生成并再次比较了三个完整的配对。并未通过放宽比较器
来使二者匹配。

这一结果表明，对于已覆盖的构造，生成的编辑器并不依赖于模型是使用 Xtext 还是
Ecore 编写的。但该结果不足以断言任意可能的 Ecore 元模型都具有等价性；要进行
这种推广，还需要增加所评估输入路线的数量和多样性。

\section{补充技术证据}
\label{sec:evidencia-complementaria}

除外部评估外，项目还保留了针对实现的测试。六组最小化且彼此互补的
Ecore--\texttt{.m2b} 配对检查了从两类输入到可执行编辑器的完整流程，并验证了
控件、连接、工具箱和生成器的多种变体。专门设计的 Ecore 案例进一步覆盖了没有文本
等价形式的构造。145 项 JUnit 测试覆盖适配器、验证和转换，而通过更新站点进行
安装则验证了插件的打包。

这些测试与语料集评估关注的问题不同。它们为所构建系统的回归、集成和部署提供了证据，
但不作为相对于现有 Blockly 编辑器有效性的主要证明。同样，应用的渲染层---例如迷宫、
画布或音乐模拟---不在主要比较范围之内：Model2Blockly 生成编辑器及其生成器，而不生成
负责执行所产出程序的领域引擎。

\section{效度威胁}
\label{sec:amenazas-validez}

\paragraph{语料集选择。}
十个案例均来自官方代码仓库，并覆盖不同的规模与机制，但它们构成的是一个目的性
样本。这些百分比描述的是该语料集，而不是所有 Blockly 编辑器的统计总体。

\paragraph{依赖项的演化。}
原始编辑器使用较早版本的 Blockly 开发。为避免后续更新改变比较结果，研究固定了
源代码修订版本和 Blockly 13.1.1 版本，并在同一个兼容性容器中运行两个处理方案。
这一结果衡量的是编辑器配置，而不是每个完整应用的历史兼容性。

\paragraph{粒度与分类。}
将评估对象分解为属性并划分为 \texttt{partial}、\texttt{unsupported} 和
\texttt{excluded} 状态，需要做出建模决策。为减少偏差，研究通过两项试点固定了
分类方案，为每项差异指定了一条具有明确标识的规则，并同时公布计数和规范描述符。即便如此，
采用另一种粒度仍可能改变绝对百分比。

\paragraph{产物的独立性。}
Pond Tutor 与 Pond 共享官方源代码。主要计算将二者视为不同的可维护配置，保守
计算则剔除重复代码行。缩减率仍处于 70,00\,\% 至 72,48\,\% 之间，由此可将
这种依赖关系的影响界定在这一范围内。

\paragraph{视觉与语义等价性。}
截图不能证明视觉效果完全相同。评估比较了结构和生成器，但没有穷举运行所有可能的
程序，也没有比较领域引擎的行为。因此，报告的一致率不应被理解为 Blockly Games
应用之间的完全等价性。

\paragraph{工作量度量。}
LOC 和字节数是可复现的源代码规模指标，并不衡量开发时间或认知难度。评估
证明了所维护规格的简洁性；若要断言时间或人力成本也有同等程度的降低，则需要开展
参与者研究。

\section{评估问题的回答}
\label{sec:respuesta-rq}

\begin{table}[htbp]
  \centering
  \caption{评估问题及回答汇总}
  \label{tab:respuesta-rq}
  \footnotesize
  \begin{tabular}{L{0.24\textwidth}L{0.68\textwidth}}
    \toprule
    评估问题 & 评估结论 \\
    \midrule
    \textbf{RQ1：功能复现} & 十个编辑器均可加载；加权结构一致率为 83,82\,\%，生成器一致率为 78,01\,\%。所有案例均被归类为部分复现。 \\
    \textbf{RQ2：规格编写工作量} & M2B 将 LOC 减少了 72,48\,\%；剔除共享的官方源代码后，保守计算的缩减率为 70,00\,\%。 \\
    \textbf{RQ3：局限} & 主要缺失能力涉及连接策略、工具箱组合、呈现用输入，以及带有状态、变形器或自动序列化机制的生成器。 \\
    \textbf{RQ4：输入路线} & 在选定的三个案例中，Ecore 与 \texttt{.m2b} 路线生成严格等价的规范描述符（3/3）。 \\
    \bottomrule
  \end{tabular}
\end{table}

\section{评估结论}
\label{sec:conclusion-evaluacion}

本次评估的证据范围比单个示例更广。对十个官方编辑器的结果表明，\allowbreak{}Model2Blockly
复现了大部分属性，并将需维护的配置源代码减少约 70\,\%。三个
Ecore--\texttt{.m2b} 配对均等价，这进一步支持了适配器与通用规格相分离的设计。
但并未实现完全等价：所有案例仍然存在差异，尤其是在以编程方式构建的工具箱、
变形器或有状态生成方面。因此，该语料集量化了当前覆盖范围、已经实现的缩减幅度，
以及扩大覆盖范围所需的扩展。
